\section{Case-study design}
\label{sec:case_studies}

Three datasets are chosen to span the assumption space, not to flatter the framework:
one where the physics is rich and documented (the debutanizer), one that supports
controlled regime structure for migration and coverage studies (TEP), and one that is
deliberately hostile --- high-dimensional, anonymized, physics-free (SECOM). The same
pipeline, validation protocol, and hyper-parameter discipline run on all three.

\hypertarget{debutanizer-column}{%
\subsection{Debutanizer column}\label{debutanizer-column}}

The benchmark of \citep{fortuna2007}: 2,394 samples of seven process inputs
(temperatures, pressure, reflux and flow signals, \(u_1\)--\(u_7\)) and the bottoms C4
content \(y\), pre-normalized to \([0,1]\). The final 360 samples are held out as the test
window; the preceding 2,034 form the selection pool for blocked CV (\(K = 5\)). The
transport lag is diagnosed at 15 samples (§3.1) and the label delay is set to the
documented analyzer delay \(\theta = 4\). This dataset carries the feature-ablation
(C2), lag-diagnosis (C3), and drift/bias-update (C4) results.

\hypertarget{tennessee-eastman-process-regimes}{%
\subsection{Tennessee Eastman process regimes}\label{tennessee-eastman-process-regimes}}

The TEP \citep{downs1993} supplies the regime structure for the coverage (C7, C8)
and migration (C6) studies. Three operating-point regimes of 2,000 rows each (100 h at
the 3-minute analyzer cadence; 41 measured + 12 manipulated variables, of which the 22
continuous measurements and the manipulated variables serve as inputs --- the delayed
composition analyzers are excluded except as the target, \(y\) = XMEAS 40, product-G) were generated with the closed-loop
Ricker/Russell--Chiang--Braatz simulation lineage \citep{bathelt2015}, with product-G
composition means of 53.8, 58.2 and 63.4 mol\% --- between-regime gaps of 4.4 and 5.2
mol\% against a within-regime standard deviation of ≈1.3, so the regimes are well
separated without being disjoint. Two framing notes are part of the design, not
caveats discovered later. First, these are \emph{feed-ratio operating points}, not the
canonical Downs--Vogel G/H product modes; the canonical modes were obtained as a
supplementary check and proved textbook-clean but under-excited for soft-sensor
training (within-mode \(R^2 \approx 0.04\)), which is precisely why deliberate
excitation was injected for the headline data. Second, an open Simulink
reimplementation \citep{vosloo2025} was abandoned for a non-reproducible
reactor-pressure startup trip; reproducibility of the data generator was treated as a
gating requirement. Splits are time-ordered 70/15/15 per regime; label delays
\(\theta \in \{2, 5\}\) bracket the plausible analyzer range; the conformal protocol is
fixed at \(\alpha = 0.10\), \(\gamma = 0.05\), window 200 (§3.4).

\hypertarget{secom-negative-control}{%
\subsection{SECOM (negative control)}\label{secom-negative-control}}

The UCI SECOM dataset \citep{mccann2008secom}: 1,567 production lots × 590 anonymized
sensor features from a semiconductor line over three months of 2008, with a pass/fail
label (104 fails, 6.6\%), 4.5\% missing cells overall, 32 features more than 40\%
missing, and 116 near-constant features --- \(p \approx n\), heavy missingness, severe
imbalance, no physics anchors. Because the framework's online machinery is regression,
SECOM is framed as \emph{virtual metrology}: one continuous measurement is held out as the
target, selected by stated, audited criteria (missingness ≤ 5\%, non-degenerate
variance, maximal \textbar point-biserial correlation\textbar{} with pass/fail --- the yield-relevant
measurement a VM sensor would exist to predict; selected \(x_{59}\), \(|r| = 0.156\)). The
fail label defines the problem and is never a model feature. Label-free screens
(missingness, near-zero variance) are applied once --- no target is consulted, so no
selection leakage is possible --- leaving 442 features; all supervised selection happens
inside CV folds via the elastic net's own shrinkage, with median imputation inside the
estimator pipeline (fold-train fit). Splits, delays, and the conformal protocol are
identical to §4.2.

\hypertarget{protocol-summary-and-metrics}{%
\subsection{Protocol summary and metrics}\label{protocol-summary-and-metrics}}

\textbf{T0 --- Protocol at a glance.}

\begin{table}[t]\centering\small
\caption{Protocol at a glance.}\label{tab:t0}
\begin{tabular}{@{}
  >{\raggedright\arraybackslash}p{(\columnwidth - 6\tabcolsep) * \real{0.2500}}
  >{\raggedright\arraybackslash}p{(\columnwidth - 6\tabcolsep) * \real{0.2500}}
  >{\raggedright\arraybackslash}p{(\columnwidth - 6\tabcolsep) * \real{0.2500}}
  >{\raggedright\arraybackslash}p{(\columnwidth - 6\tabcolsep) * \real{0.2500}}@{}}
\toprule\noalign{}
\begin{minipage}[b]{\linewidth}\raggedright
\end{minipage} & \begin{minipage}[b]{\linewidth}\raggedright
Debutanizer
\end{minipage} & \begin{minipage}[b]{\linewidth}\raggedright
TEP regimes
\end{minipage} & \begin{minipage}[b]{\linewidth}\raggedright
SECOM
\end{minipage} \\
\midrule\noalign{}

\bottomrule\noalign{}
\endlastfoot
samples × features & 2,394 × 7 & 3 × (2,000 × 53) & 1,567 × 590 \\
physics anchors & VLE bridge & process features & none \\
target & bottoms C4 & product-G comp. & \(x_{59}\) (VM) \\
transport lag & 15 (diagnosed) & diagnosed/regime & --- \\
label delay \(\theta\) & 4 (documented) & \{2, 5\} & \{2, 5\} \\
selection & blocked CV \(K{=}5\) + one-SE & idem & idem (elastic-net path) \\
split & pool 2,034 / test 360 & 70/15/15 time-ordered & 70/15/15 time-ordered \\
\end{tabular}
\end{table}

Point accuracy is reported as \(R^2\) (per fold, mean ± SE, worst fold, and held-out
test); uncertainty as empirical coverage against the 0.90 target with mean interval
width; transfer as data efficiency (§5.4); deployment as client-observed latency
percentiles. All runs are deterministic; every reported number is regenerated by a
single command per evidence file and committed with provenance stamps
(\texttt{docs/paper/evidence/}), and the full pipeline reproduced exactly across two machines.
